CRISP-DM propone una secuencia pr\'actica para proyectos de miner\'ia de datos: comprender el negocio, comprender los datos, prepararlos, modelar y evaluar los resultados antes de comunicar conclusiones \citeA{crispdm2000}. En este caso, la fase de negocio se traduce en preguntas sobre precios de vivienda, criminalidad y condiciones socioecon\'omicas.

El An\'alisis de Componentes Principales (PCA) es una t\'ecnica de reducci\'on de dimensionalidad que transforma un conjunto de variables correlacionadas en nuevas variables ortogonales llamadas componentes principales \citeA{jolliffe2016}. Su utilidad radica en condensar informaci\'on sin perder excesiva varianza y en revelar estructuras latentes entre variables. Dado que las variables del Boston Housing tienen escalas distintas, la estandarizaci\'on es obligatoria antes de aplicar PCA. Las cargas de cada componente permiten interpretar qu\'e variables dominan la estructura del fen\'omeno, mientras que los scores resumen la posici\'on de cada observaci\'on en el espacio reducido.

El paquete \texttt{MASS} provee el conjunto de datos y forma parte de la tradici\'on estad\'istica aplicada en R \citeA{masspackage,venables2002}. En consecuencia, el an\'alisis combina estad\'istica descriptiva, visualizaci\'on y modelamiento lineal para justificar cada respuesta.